% ===================================================================
%  TÁBLA Uimh. 9 — Na Sléibhte agus an Baile Spá. — An Choill agus an
%  tSealgaireacht.
%  Traduction 2026 d'apres texto/io, controlee sur texto/fr, texto/en
%  et texto/es. Consigne : voir texto/ga/10-tabla-01.tex.
%
%  LE TABLEAU DE LA MONTAGNE. Le relief est celtique de bout en bout —
%  « sliabh » la montagne, « carraig » le rocher, « gleann » la vallee,
%  « tobar » la source, « coill » le bois, « uaimh » la grotte — et
%  c'est le meme resultat qu'en roumain, en ukrainien et en basque.
%  Quatre langues sans lien entre elles.
% ===================================================================

%%K t09-apar-1 apar
\VUsaut{4.0mm}
\VUtitre{15.0pt}{60}{TÁBLA Uimh. 9}
\VUsaut{3.0mm}

%%K t09-tit-1 sub
\VUcentre{12.0pt}{0}{\textit{An Chéad Radharc.}}
\VUsaut{3.0mm}

%%K t09-tit-2 sub
\VUpk{11.4pt}{40}{Na Sléibhte. --- An Baile Spá.}
\VUsaut{3.0mm}

%%K t09-c1-01-1 p
1. --- Táim seachtain i bPratz. Ní féidir liom an t-iontas ar fad a chur in iúl a
mhothaigh mé agus mé i mo sheasamh os comhair \VUgras{shliabhraon} iontach na
bPiréiní \textsuperscript{(1)}, a ngearrtar a \VUgras{chíor}
\textsuperscript{(2)} ar an spéir ghorm mar shábh ollmhór.

%%K t09-c1-02-1 p
2. --- Níor shamhlaigh mé go bhféadfadh \VUgras{beanna} na bPiréiní
\textsuperscript{(3)} teacht ar a leithéid de \VUgras{airde}, agus d'fhan mé faoi
alltacht ag na \VUgras{hoighearshruthanna} iontacha \textsuperscript{(5)} agus ag
na \VUgras{stuaiceanna} sneachta \textsuperscript{(4)}, ar a réabann
\VUgras{maidhmeanna sneachta} chomh scanrúil sin.

%%K t09-tit-3 sub
\VUsaut{3.0mm}
\VUpk{10.2pt}{40}{Radharc sléibhe.}
\VUsaut{3.0mm}

%%K t09-c1-03-1 p
3. --- An lá tar éis teacht dom chuaigh mé chuig an sean-\VUgras{bholcán} múchta
\textsuperscript{(6)} in aice le Pratz. Ar imeall lom aimrid an \VUgras{chráitéir}
feictear fós go soiléir na rianta a d'fhág an \VUgras{laibhe}, a shil cúpla céad
bliain ó shin síos \VUgras{fána an tsléibhe} \textsuperscript{(7)}. Ar cheann de na
\VUgras{fánaí} d'éirigh liom cúpla píosa den laibhe sin a aimsiú.

%%K t09-c1-04-1 p
4. --- Tá Pratz ar an teorainn, agus tá an \VUgras{daingean}
\textsuperscript{(8)}, a chosnaíonn an \VUgras{bhearna} \textsuperscript{(27)} idir
an Fhrainc agus an Spáinn, an-chosúil leis na sean-chaisleáin ar bhruacha na
Réine, a dhealraíonn a bheith ag faire ar \VUgras{chreach} ó bharr a gcarraige.

%%K t09-c1-05-1 p
5. --- Is minic a tharraingíonn \VUgras{iolar} ollmhór \textsuperscript{(9)} m'aird,
iolar a bhfuil a nead in aice leis an \VUgras{dún} \textsuperscript{(8)}. Is minic
a thugann an \VUgras{t-éan creiche} seo, lena \VUgras{ghob} cumhachtach, uan nó
meannán chuig a scalltáin, agus é fáiscthe ina \VUgras{chrúba}. Tá a
\VUgras{sciatháin} chomh mór sin gur féidir leis ainliú go fada lena ualach trom.

%%K t09-tit-4 sub
\VUsaut{3.0mm}
\VUpk{10.2pt}{40}{An Coisí.}
\VUsaut{3.0mm}

%%K t09-c1-06-1 p
6. --- Is í an tsiúlóid is taitneamhaí anseo an turas chuig \VUgras{eas}
\textsuperscript{(10)} Coo. Tuiteann an \VUgras{t-uisce ag titim} ina chuilithe agus
imíonn sé ansin ina \VUgras{shruthán} gleoite sa \VUgras{choill ghiúise}
\textsuperscript{(11)} siar ón mbaile. D'ardaíomar tríd an gcoill sin go dtí an
\VUgras{stáisiún meitéareolaíoch} \textsuperscript{(12)} agus ó bharr an
\VUgras{túir} ghlacamar isteach lenár súile \VUgras{lánléargas} an cheantair ar
fad. Tá an \VUgras{radharc} iontach, agus dealraíonn go bhfuil an
\VUgras{dúlra} tar éis gach seod a bhí ar fáil aige a chaitheamh ar an gceantar
seo.

%%K t09-c1-07-1 p
7. --- Chun teacht anuas, thógamar an \VUgras{cáblabhealach}
\textsuperscript{(13)} agus tháinig muid anuas ag \VUgras{stáisiún} beag in aice le
\VUgras{fothracha} \textsuperscript{(14)} an tsean-\VUgras{chaisleáin fheodaigh},
ina raibh tiarnaí Phratz ina gcónaí tráth.

%%K t09-c1-08-1 p
8. --- An caisleán bocht! Cad a chaithfidh sé a cheapadh agus é ag féachaint faoi ar
an \VUgras{mbaile spá} néata \textsuperscript{(15)}, atá anois ina
\VUgras{ionad uisce} faiseanta?

%%K t09-c1-08-2 p
Tá an \VUgras{aeráid} chomh taitneamhach sin anseo, agus an \VUgras{t-uisce
mianrach} chomh leighsiúil sin, gur pléisiúr fíor an \VUgras{cúrsa leighis} a
dhéantar sa \VUgras{sanatóiriam} \textsuperscript{(17)}.

%%K t09-tit-5 sub
\VUsaut{3.0mm}
\VUpk{10.2pt}{40}{Baile spá taitneamhach.}
\VUsaut{3.0mm}

%%K t09-c1-09-1 p
9. --- Go deimhin, rinneadh gach rud chun go mbeadh an fanacht taitneamhach. Tógadh
\VUgras{tarbhealach} mór \textsuperscript{(16)} don iarnród; tá \VUgras{páirc}
\textsuperscript{(18)} an \VUgras{cheoláras}, mar a dtugtar \VUgras{ceolchoirmeacha},
leagtha amach go hálainn. Tógadh roinnt \VUgras{villí} galánta
\textsuperscript{(19)} timpeall an \VUgras{cheasaíneo} \textsuperscript{(21)}, agus
déanann an \VUgras{bóthar} atá coinnithe go han-mhaith \textsuperscript{(22)} an
taisteal thar a bheith éasca.

%%K t09-c1-10-1 p
10. --- Ní dheighleann ach \VUgras{claífort} simplí \textsuperscript{(23)} an
\VUgras{bóthar} \textsuperscript{(20)} ón \VUgras{ngleann} féin
\textsuperscript{(25)}, ar a bhfuil ina íochtar \VUgras{loch} beag gleoite
\textsuperscript{(26)}, a bheathaíonn \VUgras{sruthán} glé
\textsuperscript{(24)}.

%%K t09-c1-11-1 p
11. --- Ar an taobh eile den ghleann oibrítear \VUgras{mianach} iarainn
\textsuperscript{(28)}. Chuaigh mé cúpla uair chun féachaint ar na
\VUgras{mianadóirí} ag dul síos sa \VUgras{shaft}, agus chuaigh mé isteach sa
\VUgras{ghailearaí} féin, mar a gcaitheann siad an lá ag baint na \VUgras{mianaigh},
ina bhfuil an \VUgras{miotal}.

%%K t09-c1-12-1 p
12. --- In aice leis an mianach tógadh \VUgras{teilgcheárta} mhór, chun na saibhris
a bhaintear as a úsáid ar an láthair. Ligeann an \VUgras{foirnéis soinneáin}
\textsuperscript{(29)}, a théitear leis an \VUgras{ngual} flúirseach anseo,
\VUgras{iarann teilgthe} a fháil go tapa agus go saor.

%%K t09-c1-13-1 p
13. --- Ní i bhfad ón mianach oibrítear \VUgras{cairéal} \textsuperscript{(30)}. Ní
bhaineann an \VUgras{cairéalaí} críonna lena \VUgras{phiocóid} \VUgras{eibhear},
ná \VUgras{porfaire}, ná fiú \VUgras{gaineamhchloch}, ach \VUgras{cloch thógála}
shimplí chun tithe a thógáil.

%%K t09-c1-14-1 p
14. --- Ceann de na maisiúcháin is áille sa cheantar seo ná an \VUgras{ghiúis}
\textsuperscript{(31)}. Gach áit --- in íochtar an \VUgras{ghleanna chúng}
\textsuperscript{(32)}, ar bhruach an \VUgras{tsrutha} \textsuperscript{(33)}, in
aice leis an \VUgras{duibheagán} \textsuperscript{(34)} nó leis an aill ---
tugann a snáthaidí caola a nóta glas féin.

%%K t09-tit-6 sub
\VUsaut{3.0mm}
\VUpk{10.2pt}{40}{Siúlóidí de shiúl na gcos.}
\VUsaut{3.0mm}

%%K t09-c1-15-1 p
15. --- Bainim leas as an tsaoire chun siúlóidí fada a dhéanamh. Ligeann na
\VUgras{cosáin} \textsuperscript{(35)}, a chasann ar an sliabh, cúpla
\VUgras{ciliméadar} a shiúl, agus go minic cúpla \VUgras{léig} fiú, gan an
fad a thabhairt faoi deara.

%%K t09-c1-15-2 p
Marcálann na \VUgras{clocha míle} \textsuperscript{(36)} agus na
\VUgras{comharthaí} \textsuperscript{(37)} na cosáin, ar a gcastar go minic
\VUgras{sléibhteánach} óg \textsuperscript{(38)} le \VUgras{mála}
\textsuperscript{(40)} ar a thaobh, ag iompar \VUgras{marmait}
\textsuperscript{(39)}, nó \VUgras{giofóg} éigin \textsuperscript{(41)}, a
gcuireann a \VUgras{caipín fionnaidh} \textsuperscript{(42)} caismirt aisteach lena
\VUgras{clóca} sean stróicthe \textsuperscript{(43)}. Tugann sí ar
\VUgras{shlabhra} \VUgras{béar} foghlamtha \textsuperscript{(44)}.

%%K t09-tit-7 sub
\VUpk{10.2pt}{40}{Dreapadóireacht.}
\VUsaut{3.0mm}

%%K t09-c1-16-1 p
16. --- Beagnach gach lá téim leis an \VUgras{treoraí} \textsuperscript{(48)} ar
\VUgras{dhreapadh}, agus táim an-bhródúil nuair a bhíonn \VUgras{mála droma}
\textsuperscript{(47)} ar mo dhroim agus bata sléibhe i mo láimh, mar
\VUgras{thurasóir} \textsuperscript{(46)} fíor, agus mé ag tabhairt ruathair faoi
na \VUgras{carraigeacha} \textsuperscript{(45)}.

%%K t09-c1-17-1 p
17. --- Ó bharr an tsléibhe ní dhealraíonn na daoine ar an machaire níos mó ná
seangáin; tá an \VUgras{tréad caorach} \textsuperscript{(50)}, atá ag innilt ar an
\VUgras{bhféarach} \textsuperscript{(49)}, cosúil le bréagán páiste. Níl sa
\VUgras{ghiúis} \textsuperscript{(51)} ná fiú sa \VUgras{teachín} adhmaid
\textsuperscript{(52)} ach ponc ar éigean le feiceáil ar an nglasra.

%%K t09-c1-18-1 p
18. --- Ar na siúlóidí seo castar go minic orainn ar an \VUgras{gcosán}
\textsuperscript{(53)} an \VUgras{sealgaire fia gabhair} \textsuperscript{(54)}, ag
iompar ar a ghualainn an ainmhí a mharaigh sé ar ball.

%%K t09-c1-19-1 p
19. --- Chuir mé i mo luibhchlár duilleog an-suntasach \VUgras{raithní}
\textsuperscript{(55)} agus craobhóg bhándearg ghleoite \VUgras{fraoigh}
\textsuperscript{(57)}, a bhain mé ag béal \VUgras{uaimhe} bige
\textsuperscript{(56)} ar an tsiúlóid dheireanach.

%%K t09-tit-8 sub
\VUsaut{3.0mm}
\VUcentre{12.0pt}{0}{\textit{An Dara Radharc.}}
\VUsaut{3.0mm}

%%K t09-tit-9 sub
\VUpk{11.4pt}{40}{An Choill. --- An tSealgaireacht.}
\VUsaut{3.0mm}

%%K t09-c2-01-1 p
1. --- Inné chaith mé lá aoibhinn sa choill. Chuir beatha ghníomhach na n-ainmhithe
agus na \VUgras{bhfeithidí} \textsuperscript{(5)} a chónaíonn ann suim mhór ionam.

%%K t09-c2-01-2 p
An \VUgras{damhán alla} \textsuperscript{(1)}, ag fí a \VUgras{líon}
\textsuperscript{(2)} idir dhá chraobh chun greim a fháil ar an \VUgras{gcuileog}
\textsuperscript{(3)} atá ag dul thar bráid; an \VUgras{seilide}
\textsuperscript{(4)}, ag tarraingt go trom lena theachín; an priompallán, ag faire
ar chreach; an seangán dícheallach, ag obair in aice leis an
\VUgras{seangánlann} \textsuperscript{(6)} --- bhí na créatúir bheaga seo go léir
ag corraí agus ag saothrú le neart neamhghnách.

%%K t09-c2-02-1 p
2. --- An \VUgras{nathair} \textsuperscript{(7)}, ar shíl mé ar dtús gur
\VUgras{nathair nimhe} í, agus nach raibh inti ach \VUgras{sciathán leathair}
neamhurchóideach, chas sí í féin faoi stoc crainn agus chuir sí a ceann tanaí
amach ó am go ham, í i gcónaí mar a bheadh sí réidh le cealg. Os mo chomhair bhí
\VUgras{drúchtín} mór \textsuperscript{(8)} ag lámhacán go mall, drúchtín a d'ith
ceann de na \VUgras{sútha talún} beaga blasta \textsuperscript{(11)}, a bhfuil an
oiread sin díobh anseo.

%%K t09-c2-03-1 p
3. --- Bhí an talamh clúdaithe le \VUgras{bláthanna coille}: bhí na
\VUgras{sabhaircíní} \textsuperscript{(9)}, na \VUgras{bréanáin}
\textsuperscript{(10)} agus go leor plandaí eile nach raibh aithne agam orthu faoi
bhláth i measc na \VUgras{dtorthóg féir} \textsuperscript{(12)}.

%%K t09-c2-04-1 p
4. --- Sa cheantar seo tá an \VUgras{strufal} beagnach chomh coitianta leis an
\VUgras{muisiriún} \textsuperscript{(14)}, agus bhí \VUgras{banbh}
\textsuperscript{(13)} le maor coille éigin ag rómhar go santach, féachaint an
bhfaigheadh sé iad.

%%K t09-tit-10 sub
\VUsaut{3.0mm}
\VUpk{10.2pt}{40}{An Phóitseáil.}
\VUsaut{3.0mm}

%%K t09-c2-05-1 p
5. --- Rug mé ar \VUgras{dheireadh} radhairc a chuir an-ghreann orm. A bhuí le
boladh a \VUgras{dhaicseoinn} \textsuperscript{(15)}, bhí an \VUgras{maor coille}
\textsuperscript{(16)} díreach tar éis breith ar an \VUgras{bpóitseálaí}
\textsuperscript{(22)} agus é ar tí an \VUgras{fhiabheatha}
\textsuperscript{(17)} a mharaigh sé a bhreith leis. Ó b'fhearr leis scaradh lena
chreach ná bheith gafa, theith an póitseálaí, agus d'fhan an \VUgras{giorria}
\textsuperscript{(18)}, an \VUgras{eilit} \textsuperscript{(19)}, an
\VUgras{fia rua} \textsuperscript{(20)} agus an \VUgras{torc}
\textsuperscript{(21)} ina luí ar an bhféar, chun áthais mhóir an mhaoir choille.

%%K t09-c2-06-1 p
6. --- Cuireann éagsúlacht na gcrann sa choill seo iontas ort. Fíonn na
\VUgras{feá} \textsuperscript{(23)} agus na \VUgras{beitheanna}
\textsuperscript{(26)}, na \VUgras{giúiseanna}, a thugann \VUgras{roisín}
\textsuperscript{(27)}, na \VUgras{daracha} \textsuperscript{(29)} agus go leor
eile a gcraobhacha le chéile.

%%K t09-c2-06-2 p
Tugann an \VUgras{t-eidhneán} \textsuperscript{(24)}, a ghreamaíonn de stoic na
gcrann ard, glas geal don \VUgras{choill ard} \textsuperscript{(25)}, glas a
mheascann go taitneamhach le ton bán bog na \VUgras{caonaí}
\textsuperscript{(31)}, mar a mbíonn an \VUgras{snag darach glas}
\textsuperscript{(30)} ag cuardach feithidí.

%%K t09-tit-11 sub
\VUsaut{3.0mm}
\VUpk{10.2pt}{40}{Radharc coille.}
\VUsaut{3.0mm}

%%K t09-c2-07-1 p
7. --- Chuir na sean-daracha draíocht orm. Tugann a \VUgras{bhfréamhacha} móra
snaidhmthe \textsuperscript{(32)}, atá leath amuigh as an talamh, cuma iontach
nirt agus cumhachta dóibh, agus ní thuigim conas is féidir leis an
\VUgras{n-úinéir} \textsuperscript{(35)} cinneadh iad a leagan agus iad a thabhairt
do \VUgras{shábh} \textsuperscript{(34)} an \VUgras{tuadóra}
\textsuperscript{(33)}. Cuireann na \VUgras{stoic} bhochta \textsuperscript{(36)},
gan \VUgras{choirt} \textsuperscript{(37)} orthu nó scoilte le \VUgras{tua}
\textsuperscript{(38)} an \VUgras{fhir lae} \textsuperscript{(39)}, brón orm.

%%K t09-c2-08-1 p
8. --- Lena thaobh rinne \VUgras{gualadóir} sean \textsuperscript{(41)} lena
\VUgras{shluasaid} \VUgras{carn} \textsuperscript{(42)} guail, agus bhí seanbhean
ag iompar \VUgras{beart} \textsuperscript{(40)} \VUgras{brosna} as craobhacha na
gcrann leagtha.

%%K t09-tit-12 sub
\VUsaut{3.0mm}
\VUpk{10.2pt}{40}{An tSealgaireacht.}
\VUsaut{3.0mm}

%%K t09-c2-09-1 p
9. --- Bhí sé ar intinn agam scíth bheag a ghlacadh i \VUgras{dteach an mhaoir
choille} \textsuperscript{(46)}, ach, ag teacht chuig an \VUgras{lochán} beag
\textsuperscript{(43)}, ar a raibh \VUgras{eala} álainn \textsuperscript{(45)} ag
snámh, thug mé faoi deara go raibh an \VUgras{tuile} le déanaí
\textsuperscript{(44)} tar éis an cosán a dhéanamh ina \VUgras{riasc} ceart. Bhí orm
cor a chur díom, agus, ag dul thar an \VUgras{gcéadar} \textsuperscript{(49)}, na
\VUgras{poibleoga} \textsuperscript{(48)} agus an \VUgras{leamhán}
\textsuperscript{(47)} atá ar imeall na coille, chonaic mé conas mar a bhris as an \VUgras{mhuine}
\textsuperscript{(54)} \VUgras{fia} iontach \textsuperscript{(50)}, agus
\VUgras{conairt} dhothuirsithe \textsuperscript{(51)}
sa tóir air, á griosú ag \VUgras{adharc} \textsuperscript{(53)} na
\VUgras{sealgairí ar muin capaill} \textsuperscript{(52)}. Bhí an t-ainmhí bocht
spíonta ar fad. Thit sé, ag fáil bháis, cúpla coiscéim uaim.

%%K t09-c2-10-1 p
10. --- Tháinig mac an mhaoir choille, tarraingthe ag torann na
\VUgras{sealgaireachta}, amach as an teach, agus d'fhilleamar beirt ar an gcoill.
Thug sé faoi deara \VUgras{snag breac} \textsuperscript{(56)} ar chraobh
sean-daraigh. Cheap sé go raibh a nead i bhfolach sa \VUgras{duilliúr}
\textsuperscript{(58)}, agus theastaigh uaidh í a thógáil.

%%K t09-c2-10-2 p
Chomh haclaí le \VUgras{iora} \textsuperscript{(61)}, dhreap sé ó
\VUgras{chraobh} \textsuperscript{(55)} go craobh, ach ní bhfuair sé faic agus níor
scanraigh sé ach \VUgras{préachán} sean \textsuperscript{(60)}, a bhí ina shuí ar
\VUgras{ghéag} \textsuperscript{(57)}.
\VUsaut{4.0mm}
\VUfilet{22mm}
